% Lecture 08: Design, SOLID, and patterns
\section{Лекция 8. Design, SOLID и паттерны}
\label{sec:lecture-08}

В блоке~7 мы получили runtime polymorphism. Но корректный синтаксис ещё не
даёт удобный design: программа может компилироваться и заставлять менять один
центральный \texttt{switch} при каждом новом требовании.

\begin{importantbox}[Главная причинная цепочка]
Простой работающий design $\rightarrow$ новое требование $\rightarrow$
неудобные связи $\rightarrow$ design pressure $\rightarrow$ pattern и его
цена. Pattern --- повторяемая идея взаимодействия объектов, не готовый кусок
синтаксиса C++.
\end{importantbox}

\subsection{Design после синтаксиса}

\textbf{Responsibility} --- связная обязанность и осмысленная причина менять
модуль. \textbf{Cohesion} высока, когда его данные и операции служат этой
обязанности. \textbf{Coupling} показывает, сколько одна часть знает о другой.
\textbf{Dependency} означает, что одной части нужна другая. \textbf{Interface}
задаёт контракт, \textbf{implementation detail} можно менять за ним, а
\textbf{extension point} принимает ожидаемый новый вариант.

Например, checkout с формулами всех скидок зависит от каждой policy. Checkout,
вызывающий \texttt{PricingStrategy::price}, зависит от узкого interface;
формула стала detail. Coupling не исчез, но dependency направлен к более
стабильной abstraction. Главный тест design --- изменение требований.

\subsection{SOLID --- карта проблем, не догма}

SOLID --- guidelines design, не требования стандарта C++ и не законы языка.
Compiler не способен полностью проверить их смысл.

\paragraph{SRP.} Single Responsibility Principle: class/module имеет одну
связную responsibility, один осмысленный reason to change. Это не «один method
на class». Несколько измерений \texttt{CourseBlock} связны; чтение filesystem,
выбор policy и три формата вывода имеют разные причины изменения. Разделение
повышает cohesion, но лишние классы для одноразовой функции только мешают.

\paragraph{OCP.} Open/Closed Principle: стабильный central code желательно
расширять вариантами без постоянного редактирования. Абсолютной закрытости нет:
интерфейс и требования меняются. Extension point создают лишь на ожидаемой
оси вариативности; иначе это speculative complexity.

\paragraph{LSP.} Liskov Substitution Principle связывает design с блоком~7:
derived сохраняет base contract. Если pricing interface обещает
неотрицательную цену, override с \texttt{-1} ломает ожидание client, хотя
signature корректна. Compiler проверяет \texttt{override}, но не поведенческий
смысл. Client не должен делать исключения по concrete subtype.

\paragraph{ISP.} Interface Segregation Principle: client не зависит от
операций большого interface, которые ему не нужны. Result observer требует
\texttt{on\_result}, но не \texttt{save\_html} и \texttt{set\_color}. Маленькие
целевые interfaces полезны при разных clients; дробить interface без причины
тоже не нужно.

\paragraph{DIP.} Dependency Inversion Principle: high-level policy не
привязана жёстко к low-level implementation; обе зависят от abstraction.
Простейшая dependency injection --- constructor parameter:
\begin{cppcode}
explicit Checkout(const PricingStrategy& strategy) : strategy_{strategy} {}
\end{cppcode}
DI container/framework не требуется. Reference non-owning: strategy должна
жить дольше checkout.

\subsection{Strategy: независимо заменяемая policy}

Baseline с двумя стабильными вариантами честен:
\begin{cppcode}
int direct_price(int base, PricePolicy policy) {
    switch (policy) {
    case PricePolicy::regular: return base;
    case PricePolicy::student: return base * 80 / 100;
    }
    return base;
}
\end{cppcode}
Pressure появляется, когда algorithms добавляются независимо, выбираются
runtime, а central switch приходится править для каждого варианта.

\begin{cppcode}
class PricingStrategy {
public:
    virtual ~PricingStrategy() = default;
    virtual int price(int base) const = 0;
};
class StudentPricing final : public PricingStrategy {
public:
    int price(int base) const override { return base * 80 / 100; }
};
\end{cppcode}

\texttt{PricingStrategy} --- Strategy, конкретные тарифы ---
ConcreteStrategy, \texttt{Checkout} --- Context. Caller выбирает object, а
стабильный Context знает только abstraction. Compiler проверяет hierarchy,
pure virtual contract и \texttt{override}; runtime virtual dispatch выбирает
реализацию по dynamic type. Templates будут в блоке~9 и здесь не являются
основным механизмом.

В лаборатории automatic strategy живёт дольше Context; \texttt{const\&} не
владеет ею. Цена: extra object/type, indirection и обычно один virtual call.
Allocation не обязательна. Для короткого неизменного switch Strategy хуже.
В hot path стоимость измеряют; вне него design value часто важнее одного call.

\subsection{Factory Method: выбор product у creator hierarchy}

Constructor, helper \texttt{make\_x} и simple factory с одним switch могут
быть хороши, но это не GoF Factory Method. Здесь pressure иной: общий алгоритм
Creator работает через Product interface, а ConcreteCreator решает, какой
ConcreteProduct предоставить.

\begin{cppcode}
class ReportCreator {
public:
    std::string publish(const std::string& title) const {
        const Product& product = create_product();
        return title + " -> " + product.format(title);
    }
protected:
    virtual const Product& create_product() const = 0;
};
\end{cppcode}

\texttt{Product} --- interface; \texttt{TextProduct}/\texttt{HtmlProduct} ---
ConcreteProduct; \texttt{ReportCreator} --- Creator с factory method
\texttt{create\_product}; два derived creators --- ConcreteCreator. Именно
creator polymorphism, а не слово \texttt{create}, делает пример Factory Method.

Каждый ConcreteCreator владеет product как member и возвращает
\texttt{const Product\&}. Reference non-owning и действительна лишь пока жив
creator. Heap allocation не нужна. Production factory часто возвращает owning
smart pointer, но ownership tools относятся к блоку~11. Цена --- две
hierarchies, virtual calls, types и lifetime contract. Для одного product
constructor/helper проще.

\subsection{Observer: publisher не знает всех получателей}

Baseline publisher напрямую вызывает A, B и C. Новый receiver требует менять
publisher. Pressure: несколько независимо добавляемых consumers реагируют на
одно event.
\begin{cppcode}
class Observer {
public:
    virtual ~Observer() = default;
    virtual void on_completed(const std::string& job, int result) = 0;
};
void subscribe(Observer& observer) { observers_.push_back(&observer); }
std::vector<Observer*> observers_; // non-owning
\end{cppcode}

Subject/Publisher хранит registry abstractions; ConcreteObserver реализует
effect; event --- завершение job; \texttt{notify} вызывает каждого receiver.
Automatic observers обязаны жить дольше регистрации и notify. Subject ничего
не удаляет. Уничтоженный observer оставит dangling pointer и возможное UB.
В долгоживущей системе нужен unsubscribe/registration design. Smart pointers
не универсальное лекарство: shared ownership тоже имеет цену и cycles.

Notify стоит $O(n)$ virtual calls и storage для $n$ observers; появляются
вопросы порядка callbacks и reentrancy. Для одного постоянного receiver прямой
call честнее.

\subsection{Pattern, feature, library и композиция}

\texttt{virtual}, inheritance, pointer и reference --- language features.
Strategy, Factory Method и Observer --- patterns, реализуемые этими features.
Позже возможны templates, lambdas, \texttt{std::function} и smart pointers;
реализация не меняет design pressure. Patterns можно сочетать, но не обязаны:
Strategy выбирает policy, Observer сообщает результат, Factory Method иногда
делегирует product creator. «Паттерн-комбайн» без потребности не нужен.

\subsection{Эксперимент 1: patterns\_lab}

\texttt{01\_patterns\_lab} проверяет direct switch и две strategies с общим
workflow и разными totals; общий \texttt{Creator::publish} через два creators
и products; два events для console/statistics observers. Контрольный case
оставляет стабильный двухвариантный выбор обычным conditional.
\begin{shellcode}
cd examples/08_design_solid_and_patterns/01_patterns_lab
cmake -S . -B build -DCMAKE_BUILD_TYPE=Debug
cmake --build build
./build/patterns_lab
ctest --test-dir build --output-on-failure
\end{shellcode}
PASS: отдельные markers \texttt{STRATEGY\_CASE},
\texttt{FACTORY\_METHOD\_CASE}, \texttt{OBSERVER\_CASE},
\texttt{SIMPLE\_CASE} и итоговый \texttt{PATTERNS\_LAB=PASS}; exit code~0.

\subsection{Эксперимент 2: course\_audit\_patterns}

Read-only utility исследует реальные 15 pairs. \texttt{AuditEngine} зависит
от \texttt{AuditStrategy}. Structural проверяет наличие pair, Strict --- ещё
непустую лекцию и examples. Третья, lecture-only policy, локально проверяет
расширяемость: она добавлена без правки engine workflow.

Три observers получают каждый result: console summary, issue detail и
statistics. Engine не знает concrete effects. Strategies и observers созданы
раньше engine, поэтому его non-owning references/pointers действительны.
\begin{shellcode}
cd examples/08_design_solid_and_patterns/02_course_audit_patterns
cmake -S . -B build -DCMAKE_BUILD_TYPE=Debug
cmake --build build
./build/course_audit_patterns --root ../../.. --policy structural
./build/course_audit_patterns --root ../../.. --policy strict
ctest --test-dir build --output-on-failure
\end{shellcode}
Structural ожидает 15 pairs. Strict честно находит готовые 01--08 и семь
будущих incomplete blocks. PASS означает 15 доставленных events и ожидаемый
итог policy, а не готовность блоков 09--15.

Factory Method здесь сознательно не применён: CLI choice короток, objects
локальны, общего polymorphic Creator algorithm нет. Добавление creator/product
hierarchies было бы overengineering.

\subsection{SOLID-аудит практики}
SRP: inspection, evaluation и reporting разделены; \texttt{main} остаётся уместным composition root; OCP: strategy и observer добавляются к неизменному engine; CLI mapping для нового пользовательского имени всё же меняется; LSP: каждая policy выдаёт осмысленный result для любого block, каждый observer принимает любое событие без проверки concrete publisher; ISP: два interfaces узки под настоящих clients; дальше дробить незачем; DIP: high-level engine зависит от abstractions, injection выполнена reference parameter и \texttt{subscribe}, без framework.

\subsection{Overengineering и cost model}

Pattern добавляет objects, types, pointer/reference indirection, virtual
dispatch, code size и compile complexity. Factory Method не требует heap;
Observer notify имеет $O(n)$; Strategy runtime обычно даёт одну дополнительную
indirection. Обычная реализация virtual использует vtable/vptr, но стандарт
не требует такого layout. Performance зависит от implementation и hot path:
неверны и «паттерны медленные», и «нулевая стоимость» без контекста.

Если два варианта не меняются, switch может быть яснее. Abstraction без
variation point --- unnecessary complexity. Хороший design минимален для
текущих и ожидаемых требований; pattern не самоцель.

\subsection{Проверка понимания}
После блока нужно уметь объяснить все пять SOLID как guidelines; начать с
baseline и назвать pressure; отличить Factory Method от helper; назвать роли
трёх patterns; доказать non-owning lifetime; отделить pattern от feature;
оценить trade-offs и отказаться от pattern, когда direct code проще.

\subsection{Источники для сверки}
Virtual calls, abstract classes, references и lifetime сверены с C++17 draft
N4659 (clauses~6, 10, 12, 15). OOP background сопоставлен с MIT 6.096
Lecture~7 и лекциями~5--7. Strategy, Factory Method и Observer употреблены в
каноническом GoF смысле; cost context --- с локальным C++ Performance TR,
target-based сборка --- с Modern CMake.
